\subsection{Other languages}

The optional argument is the Pygments lexer, so any language Pygments knows is available.
The signature is the slides project's on purpose: a fragment written for one template reads the same in the other.

\begin{codebox}[c]{sum_array.c}
size_t sum_array(const int *values, size_t n) {
    size_t total = 0;
    for (size_t i = 0; i < n; i++) {
        total += values[i];  /* no overflow check, on purpose */
    }
    return total;
}
\end{codebox}

\begin{codebox}[haskell]{Fib.hs}
fib :: Int -> Integer
fib n = fibs !! n
  where fibs = 0 : 1 : zipWith (+) fibs (tail fibs)
\end{codebox}

\subsection{Without a title bar}

\code{codeplain} drops the tab and the line numbers, for an excerpt short enough that neither earns its space.

\begin{codeplain}[sql]
SELECT course, AVG(grade) AS mean
  FROM submissions
 GROUP BY course
HAVING COUNT(*) > 5;
\end{codeplain}

\subsection{Inline}

Set an identifier with \code{\textbackslash code}, as in \code{binary\_search(values, target)}, and a path with \code{\textbackslash filepath}, as in \filepath{fragments/problems/010-example.tex}.
Both stay in Ubuntu Mono, so an identifier reads the same inline as it does inside a plate.
